\chapter{The denominators}\label{chap:denominators}

Before it sieves the numerators of a denominator~$b$ (Chapter~\ref{chap:sieve}), \ratpoints{} decides
which denominators it sieves at all, in which order, and which numerators modulo~$64$ each class of
the denominators modulo~$64$ admits. This chapter describes these choices and proves that none of
them loses a point. They are made for the polynomial searched, of degree~$d$ with the coefficients
$c_0, \dots, c_d$, that the set-up hands on (Chapter~\ref{chap:setup}), with its binary form
$F(x, z) = \sum_k c_k x^k z^{D - k}$ of the even degree $D = d + (d \bmod 2)$
(Definition~\ref{def:binaryForm}). The lemmas speak of triples $(a, y, b)$ of integers~$a$, $y$
and a natural number~$b$ with $\gcd(a, b) = 1$ and $y^2 = F(a, b)$, as every point
(Definition~\ref{def:IsPoint}) is one, and the denominators range from $b_{\mathrm{low}}$ to
$b_{\mathrm{high}}$, which are~$1$ and the height bound~$H$ unless the options narrow the range
(Definition~\ref{def:regionSpec}).

The 2-adic information says which numerators modulo~$64$ each class of $b \bmod 64$ admits
(Section~\ref{sec:twoAdic}). For odd~$d$ the denominator of a point is a squarefree divisor
of~$c_d$ times a square (Section~\ref{sec:oddDegree}). For even~$d$ \ratpoints{} looks for primes that
divide no denominator of a point, or only to certain powers (Section~\ref{sec:forbidden}), and tests
the denominators by a Jacobi symbol (Section~\ref{sec:jacobi}). The loop over the denominators takes
one of four shapes that combine these tests (Section~\ref{sec:shapes}). Each test removes only
candidates that carry no point; what the tests change is the order of the points, and
Section~\ref{sec:searchDens} proves that the model visits the denominators in the order of the
specification (Definition~\ref{def:orderedPoints}). The forbidden divisors of the primes examined
for the sieve come out of the loop of the tables of the choice of the moduli, which also cuts the
loop over the denominators into blocks (Chapter~\ref{chap:moduli}); the class of~$b$ gives the sieve
its first pattern and the packing of its numerators (Chapter~\ref{chap:sieve}).

In the C code, \verb|find_points_work_1| sets the flags of the tests
(\csource{find_points.c}{4607}{4607}, \csource{find_points.c}{4905}{4943}): for even~$d$ the
checked loop and, when~$c_d$ is not a square and the caller does not forbid it (the
option~\verb|-j| of the program), the Jacobi test; for odd~$d$ the squares when $|c_d| = 1$ and the
divisors of \verb|setup_us1| otherwise. The model keeps the order of the C code
(\lname{Ratpoints.searchPoints}, \lname{Ratpoints.sievingInfo}, \lname{Ratpoints.Sieve.sieveRun}).

\section{The 2-adic information}\label{sec:twoAdic}

Write $f(x) = F(x, 1)$ and $\tilde f(t) = F(1, t)$, the reversed polynomial. In a pair $(a, b)$
with $\gcd(a, b) = 1$, $a$ or~$b$ is odd. For odd~$b$, $F(a, b) \equiv b^D f(a b^{-1}) \bmod 64$,
where~$b^D$ is the square of a unit since~$D$ is even; for even~$b$, $a$ is odd and
$F(a, b) \equiv a^D \tilde f(b a^{-1}) \bmod 64$. So two words of~$64$ bits decide which pairs of
residues modulo~$64$ can carry a point: the residues~$k$ at which $f(k)$ is a square modulo~$64$,
and the even residues~$t$ at which $\tilde f(t)$ is one.

\begin{definition}[The 2-adic information]\label{def:twoAdic}
  \uses{def:bitArrays}
  \lean{Ratpoints.Sieve.fsq, Ratpoints.Sieve.gsq, Ratpoints.Sieve.twoAdic, Ratpoints.Sieve.denBits, Ratpoints.Sieve.twoAdic_spec, Ratpoints.Sieve.getLsbD_run_twoAdic, Ratpoints.Sieve.getLsbD_fsq_coeffsMod64, Ratpoints.Sieve.getLsbD_gsq_coeffsMod64}
  \leanok
  The words~$\mathrm{fsq}$ and~$\mathrm{gsq}$, by Horner's rule on~$64$-bit words
  (\csource{find_points.c}{1006}{1014}, \csource{find_points.c}{1015}{1021}): bit~$k$
  of~$\mathrm{fsq}$ is set exactly when $k < 64$ and $\sum_{i \le d} c_i k^i$ is a square
  in~$\mathbb{Z}/64$, bit~$t$ of~$\mathrm{gsq}$ exactly when $t < 64$ is even and
  $\sum_{i \le d} c_i t^{D - i}$ is a square in~$\mathbb{Z}/64$. The words~$\mathrm{nb}[b]$ for
  $b < 64$, scattered from these (\csource{find_points.c}{983}{994},
  \csource{find_points.c}{1023}{1050}): bit~$a$ of~$\mathrm{nb}[b]$ is set exactly when~$b$ is odd
  and $a = k b \bmod 64$ for a set bit~$k$ of~$\mathrm{fsq}$, or $a < 64$ is odd and
  $b = t a \bmod 64$ for a set bit~$t$ of~$\mathrm{gsq}$. And~$\mathrm{den\_bits}$, with bit~$b$
  set exactly when $\mathrm{nb}[b] \ne 0$. These descriptions are the contracts.
\end{definition}

By the congruences above, bit~$a$ of~$\mathrm{nb}[b]$ is set exactly when~$a$ or~$b$ is odd and
$F(a, b)$ is a square modulo~$64$.

\begin{lemma}[The class of a point]\label{lem:twoAdic_point}
  \uses{def:twoAdic, def:binaryForm, def:IsPoint, def:bruteForce}
  \lean{Ratpoints.Sieve.classBit_twoAdic_of_sq_eq, Ratpoints.Sieve.getElem!_twoAdic_ne_zero, Ratpoints.Sieve.not_isPoint_of_denBits_eq_zero}
  \leanok
  Let~$a$ and~$y$ be integers and~$b$ a natural number with $\gcd(a, b) = 1$ and $y^2 = F(a, b)$.
  Then bit $a \bmod 64$ of $\mathrm{nb}[b \bmod 64]$ is set. So the word of the class of a
  denominator over which the brute-force reference finds a point is not~$0$; and if
  $\mathrm{den\_bits} = 0$, the curve has no point of any height, the points at infinity included.
\end{lemma}

\begin{proof}
  \uses{def:twoAdic, lem:sieveEntry_sound}
  \leanok
  For odd~$b$, $f(k)$ is a square modulo~$64$ at $k = a b^{-1} \bmod 64$, and the scatter sets bit
  $k b \equiv a$ of $\mathrm{nb}[b \bmod 64]$; for even~$b$, $\tilde f(t)$ is a square at the even
  $t = b a^{-1} \bmod 64$, and the scatter sets bit $a \bmod 64$ in the word of $t a \equiv b$
  (\lname{Ratpoints.isSquare\_sum\_mul\_pow\_sub\_of\_sq\_eq\_binaryForm} and its companion, in any
  commutative ring); Horner's rule on~$64$-bit words is exact modulo~$64$. A point at infinity is a
  triple $(1, y, 0)$ and also sets a bit of~$\mathrm{den\_bits}$.
\end{proof}

So the 2-adic exit, which ends the search when $\mathrm{den\_bits} = 0$ and there is no point at
infinity, loses nothing, and its condition on the points at infinity always holds. Since bit~$0$
of~$\mathrm{fsq}$ says that~$c_0$ is a square modulo~$64$ and bit~$0$ of~$\mathrm{gsq}$ that~$c_D$
is one, the exit is taken only for even~$d$ with~$c_0$ and~$c_d$ non-squares modulo~$64$.

The classes of the numerators (\lname{Ratpoints.Sieve.numClasses}, Definition~\ref{def:packing})
have the pattern~$0$ exactly when their word is~$0$. So the class test of the loop over the
denominators, that the pattern of the class of $b \bmod 64$ is not~$0$, removes no denominator of a
point, and the pattern, with which the sieve of~$b$ starts, removes no numerator of a point.

\section{Odd degree: squares and divisors}\label{sec:oddDegree}

For odd~$d$, $F(a, b) = b\, G(a, b)$ with $G(a, b) = c_d a^d + b \sum_{k < d} c_k a^k b^{d-1-k}$.
Let $(a, y, b)$ be a triple as above and~$p$ a prime with $p \nmid c_d$ and $v_p(b) = m \ge 1$.
Then $p \nmid a$, the term $c_d a^d b$ of $F(a, b)$ has the valuation~$m$ at~$p$, and every other
term has at least~$2m$. When in a sum $y^2 = \sum_i f_i$ of integers one term $p^v u$ with
$p \nmid u$ has a lower valuation than all others, $v$ is even and~$u$ is a square modulo~$p$
(\lname{Int.even\_and\_exists\_isSquare\_add\_of\_sq\_eq\_sum}); so~$m$ is even. Hence
$b = s k^2$ with~$s$ squarefree and $s \mid c_d$, and~$b$ is a square when $|c_d| = 1$. For a prime
$p \mid c_d$ there is a second test: when $v_p(b)$ is large enough, the term $c_d a^d$ alone has the
least valuation in $G(a, b)$, so that $v_p(b) + v_p(c_d) = v_p(y^2)$ is even.

\begin{definition}[The valuation]\label{def:valuation}
  \lean{Ratpoints.veryBig, Ratpoints.Setup.valuation, Ratpoints.Setup.valuation_spec, Ratpoints.Setup.fst_run_valuation}
  \leanok
  The valuation~$v$ of an integer $n \ne 0$ at~$p$ with the unit part $(|n| / p^v) \bmod p$, by
  repeated division with remainder, and $(\mathrm{veryBig}, 0)$ for $n = 0$, with
  $\mathrm{veryBig} = 1000$ (\csource{find_points.c}{422}{422}): one function for \verb|valuation|
  on GMP integers and \verb|valuation1| on \verb|long|s (\csource{find_points.c}{424}{454}). The
  contract: for $p > 1$ the loop ends, and for $n \ne 0$, $v$ is the~$p$-adic valuation of~$|n|$
  and the unit part is not~$0$.
\end{definition}

\begin{definition}[The divisors of the leading coefficient]\label{def:setupUs1}
  \uses{def:valuation, def:orderedPoints, def:maxDegree, def:tableOfPrime}
  \lean{Ratpoints.Setup.factorLcf, Ratpoints.Setup.squarefreeDivisors, Ratpoints.Setup.slope, Ratpoints.Setup.DenInfo, Ratpoints.Setup.setupUs1, Ratpoints.Setup.factorLcf_spec, Ratpoints.Setup.slope_spec, Ratpoints.Setup.setupUs1_spec, Ratpoints.Setup.map_snd_run_setupUs1}
  \leanok
  For odd~$d$ and $|c_d| \ne 1$, as \verb|setup_us1| computes them: the factorization of~$|c_d|$ by
  trial division over the primes of the table, a cofactor below the square of the next prime taken
  as a prime (\csource{find_points.c}{777}{860}); the squarefree divisors at
  most~$b_{\mathrm{high}}$ by a doubling loop, at most $\mathrm{maxDivisors} = 512$ of them
  (\csource{find_points.c}{864}{880}); and for each prime~$p$ of~$c_d$ a record of~$p$, $v_p(c_d)$
  and the slope, the least $n \ge 1$ with $(d - k)\, n > v_p(c_d) - v_p(c_k)$ for all $k < d$
  (\csource{find_points.c}{894}{912}). The records and the divisors are returned exactly when the
  factorization ends with the cofactor~$1$ and there are fewer than~$\mathrm{maxDivisors}$ divisors
  (\csource{find_points.c}{914}{935}). The contracts say so, and: for $c_d \ne 0$ the factorization
  ends with the cofactor~$1$ exactly when \ratpoints{} factors~$|c_d|$ completely in the sense of the
  specification ($\mathrm{FactorsByTrialDivision}$, Definition~\ref{def:orderedPoints}), and then
  lists the primes of~$|c_d|$ in increasing order with their valuations; the slope satisfies
  $n \ge 1$ and $1 + v_p(c_d) - v_p(c_k) \le (d - k)\, n$ for all $k < d$, where
  $v_p(0) = \mathrm{veryBig}$; for $c_d \ne 0$ and $b_{\mathrm{high}} \ge 1$ the set-up returns
  divisors exactly when~$|c_d|$ is factored completely and has fewer than~$\mathrm{maxDivisors}$
  squarefree divisors at most~$b_{\mathrm{high}}$, and these are then
  $\mathrm{binaryDivisors}(|c_d|, b_{\mathrm{high}})$, in its order.
\end{definition}

The contract of the factorization rests on the invariant
\lname{Ratpoints.Setup.FactorLcfInvariant} of its loop. By
\lname{Ratpoints.Setup.foldl\_extendDivisors\_eq\_filter}, the doubling loop lists the products
at most~$b_{\mathrm{high}}$ of the sublists of the primes, in the order of binary counting, the
first prime the lowest bit, or a prefix of them.

\begin{lemma}[The denominators of odd degree]\label{lem:odd_denominators}
  \uses{def:setupUs1, def:valuation, def:binaryForm}
  \lean{Ratpoints.isSquare_of_sq_eq_binaryForm_of_odd, Ratpoints.exists_eq_mul_mul_of_sq_eq_binaryForm_of_odd, Ratpoints.even_add_of_sq_eq_binaryForm_of_odd, Ratpoints.Setup.DenInfo.excludes_eq_false_of_sq_eq, Ratpoints.Setup.mem_iff_of_setupUs1_eq_some, Ratpoints.Setup.nodup_of_setupUs1_eq_some}
  \leanok
  Let~$d$ be odd, $a$ and~$y$ integers and~$b$ a natural number with $\gcd(a, b) = 1$ and
  $y^2 = F(a, b)$.
  \begin{enumerate}
    \item If $|c_d| = 1$, then~$b$ is a square.
    \item $b = s k^2$ for a squarefree~$s$ that divides~$c_d$.
    \item If~$p$ is a prime with $v_p(b) \ge 1$ and $v_p(c_d) + 1 \le v_p(c_k) + (d - k)\, v_p(b)$
      for every $k < d$ with $c_k \ne 0$, then $v_p(c_d) + v_p(b)$ is even. So for $b \ge 1$ the
      record of any prime~$p$ with its slope does not exclude~$b$, where a record excludes~$b$ when
      $v_p(b)$ is at least the slope and $v_p(b) + v_p(c_d)$ is odd.
  \end{enumerate}
  When the set-up returns divisors, for $c_d \ne 0$ and $b_{\mathrm{high}} \ge 1$, they are the
  squarefree divisors of~$c_d$ at most~$b_{\mathrm{high}}$, each once.
\end{lemma}

\begin{proof}
  \uses{def:setupUs1, def:valuation, def:orderedPoints}
  \leanok
  Parts (1) and~(2) are the argument at the beginning of this section. For~(3), $p \nmid a$, and in
  $G(a, b)$ the term $c_d a^d$ has the valuation $v_p(c_d)$, every other nonzero term a larger one,
  so that $v_p(b) + v_p(c_d) = v_p(y^2)$ is even; when $v_p(b)$ is at least the slope, the contract
  of the slope gives the hypothesis. The divisors
  are $\mathrm{binaryDivisors}(|c_d|, b_{\mathrm{high}})$, each once
  (\lname{Ratpoints.mem\_binaryDivisors\_iff}, \lname{Ratpoints.nodup\_binaryDivisors}).
\end{proof}

\section{Even degree: the forbidden divisors}\label{sec:forbidden}

For even~$d$, a prime~$p$ and a triple $(a, y, b)$ as above with $v_p(b) = m \ge 1$, $p \nmid a$,
and the term~$j$ of $F(a, b) = \sum_{j \le d} c_{d-j}\, a^{d-j} b^j$ has the valuation $w_j + j m$,
where $w_j = v_p(c_{d-j})$. If one term~$j_0$ alone has the least valuation, that valuation is even,
and for even~$j_0$ the unit part~$r_{j_0}$ of~$c_{d-j_0}$ is a square modulo~$p$: the argument of
the Newton polygon. A valuation~$m$ at which this fails is one that no denominator of a point has
at~$p$. For $p \nmid c_d$ and $p \mid b$, further, $y^2 \equiv c_d a^d \bmod p$ with~$a^d$ the
square of a unit; so when~$c_d$ is not a square modulo~$p$, $p$ divides no denominator of a point.
Such primes are the forbidden divisors, of two kinds: primes tested on words of~$64$ denominators at
once (\lname{Ratpoints.Sieve.ForbiddenEntry}), and primes with a mask of excluded valuations
(\lname{Ratpoints.Sieve.ForbiddenVal}).

\begin{definition}[The forbidden divisors]\label{def:forbidden}
  \uses{def:valuation, def:reciprocal, def:reversal, def:tableOfPrime, lem:primeRows}
  \lean{Ratpoints.Sieve.forbiddenAt, Ratpoints.Sieve.forbiddenValuations, Ratpoints.Sieve.forbiddenOfPrime, Ratpoints.Sieve.forbidden, Ratpoints.Sieve.forbiddenAt_spec, Ratpoints.Sieve.forbiddenValuations_spec, Ratpoints.Sieve.IsForbiddenEntry, Ratpoints.Sieve.IsForbiddenVal, Ratpoints.Sieve.forbidden_spec}
  \leanok
  Let~$d$ be even. The test of a valuation~$m$ at a prime~$p$ is true when, for the valuations~$w_j$
  and the unit parts~$r_j$ modulo~$p$, signs included, of the nonzero~$c_{d-j}$, the least
  $w_j + j m$ is taken at one~$j_0$ only and is odd, or~$j_0$ is even and~$r_{j_0}$ is not a square
  modulo~$p$ (\csource{find_points.c}{2218}{2228}); the mask of the forbidden valuations at~$p$, a
  natural number, has bit~$m$ set for the $m \ge 1$ with $p^m \le b_{\mathrm{high}}$ that the test
  excludes (\csource{find_points.c}{2229}{2234}). The forbidden divisor of an examined prime~$p$
  (\csource{find_points.c}{2642}{2676}): for $p \mid c_d$ and a nonzero mask, $p$ with its words
  when the mask holds every $m \ge 1$ with $p^m \le b_{\mathrm{high}}$, and~$p$ with its mask
  otherwise; for $p \nmid c_d$, $p$ with its words when~$c_d$ is not a square modulo~$p$, as the
  entry at infinity of its table says. Then, unless~$c_d$ is a square, the further primes~$p$ of the
  table, from the first one not examined, with $p^2 \le b_{\mathrm{high}}$ and~$c_d$ not a square
  modulo~$p$, with their words, while fewer than~$\mathrm{maxForbidden}$ forbidden divisors are
  found (\csource{find_points.c}{2709}{2772}). The contracts: a true test means that one term of
  least valuation excludes~$m$, and the mask is as described. A forbidden divisor is
  \emph{admissible} when it is the forbidden divisor of a prime of the table, made from the residues
  of the coefficients and the entry at infinity of its table, or, with words, a prime of the table
  with its words and its reciprocal at which~$c_d$ is not a square; from admissible entries the
  further primes add only admissible ones, whatever the cap.
\end{definition}

The words of a prime~$p$ are those of \lname{Ratpoints.Sieve.forbiddenRow}, $p$~words with bit~$j$
of word~$r$ clear exactly when $p \mid 64 r + j$, read from the generated pattern \verb|sieves0| for
the first~$\mathrm{numPrimes}$ primes of the table and made as \verb|sieving_info| makes them for
the others (\lname{Ratpoints.Sieve.extract\_sieves0}). For even~$d$, and while fewer
than~$\mathrm{maxForbidden}$ forbidden divisors have been found, the loop of the tables records the
forbidden divisor of each examined prime (\csource{find_points.c}{2624}{2626}); these are admissible
(\lname{Ratpoints.Sieve.isForbiddenEntry\_of\_run\_examinePrime?\_eq\_some} and its companion), so
that the choice of the moduli hands on admissible forbidden divisors only
(\lname{Ratpoints.IsSievingInfo}, Chapter~\ref{chap:moduli}).

\begin{lemma}[The forbidden divisors keep every point]\label{lem:forbidden}
  \uses{def:forbidden, def:valuation, def:binaryForm}
  \lean{Ratpoints.even_and_isSquare_of_sq_eq_binaryForm_of_even, Ratpoints.Sieve.testBit_run_forbiddenValuations_eq_false_of_sq_eq, Ratpoints.Sieve.IsForbiddenEntry.words_eq, Ratpoints.Sieve.IsForbiddenEntry.not_dvd, Ratpoints.Sieve.IsForbiddenVal.excludes_eq_false}
  \leanok
  Let~$d$ be even, $a$ and~$y$ integers and~$b$ a natural number with $\gcd(a, b) = 1$ and
  $y^2 = F(a, b)$. For every prime~$p$ and every~$b_{\mathrm{high}}$, bit~$v_p(b)$ of the mask of
  the forbidden valuations at~$p$ is clear. An admissible forbidden divisor with words has the words
  of \lname{Ratpoints.Sieve.forbiddenRow}, and for $1 \le b \le b_{\mathrm{high}}$ its prime does
  not divide~$b$; and for $b \ge 1$ no admissible forbidden divisor with a mask excludes~$b$.
\end{lemma}

\begin{proof}
  \uses{def:forbidden, def:valuation, def:reciprocal, def:tableOfPrime, lem:exitByPrime, lem:primeRows}
  \leanok
  The first statement is the Newton polygon at~$p$
  (\lname{Ratpoints.even\_and\_isSquare\_of\_sq\_eq\_binaryForm\_of\_even}), from the statement on
  a sum with one term of least valuation; bit~$0$ is never set. A prime with words from a mask has
  every valuation up to~$b_{\mathrm{high}}$ excluded; one with $p \nmid c_d$ divides no denominator
  of a point by the congruence above, the entry at infinity of its table saying, for even~$d$,
  whether~$c_d$ is a square modulo~$p$.
\end{proof}

\section{The Jacobi test}\label{sec:jacobi}

For even~$d$, an odd prime $q \mid b$ with $q \nmid c_d$ makes~$c_d$ a square modulo~$q$, by the
congruence of Section~\ref{sec:forbidden}. So for the largest divisor~$b^*$ of~$b$ prime to
$2 c_d$, the Jacobi symbol $J(c_d \mid b^*)$ is~$1$: the test rejects every denominator with an odd
number of prime factors, counted with multiplicity, at which~$c_d$ is not a square. The C code
computes the symbol by a binary algorithm or, when the leading coefficient allows it, by tables of
the non-squares modulo its odd primes and quadratic reciprocity.

\begin{definition}[The Jacobi tests]\label{def:jacobi}
  \uses{def:valuation, def:reciprocal, def:tableOfPrime, def:bitArrays, lem:reciprocal}
  \lean{Ratpoints.Sieve.jacobi, Ratpoints.Sieve.jacobi_spec, Ratpoints.Sieve.JacobiInfo, Ratpoints.Sieve.jacobiSetup, Ratpoints.Sieve.jacobiTest, Ratpoints.Sieve.IsJacobiInfo, Ratpoints.Sieve.jacobiSetup_spec, Ratpoints.Sieve.jacobiTest_spec}
  \leanok
  The binary algorithm (\csource{find_points.c}{498}{566}, \csource{find_points.c}{568}{641}) for
  $b \ge 1$ and an integer~$\ell$: subtraction and halving on $\ell \bmod b$ and~$b$, with the signs
  of $J(2 \mid m)$ and of the reciprocity, the gcd divided out of~$b$ when it is not~$1$. Its
  contract: the result is $J(\ell \mid b^*)$, where~$b^*$ is the largest divisor of~$b$ prime
  to~$2\ell$. The tables (\csource{find_points.c}{683}{708}), for $\ell = \pm 2^v \prod_i q_i^{e_i}$
  and $b_{\mathrm{high}} \le \mathrm{mulModLimit} = 2^{32} - 1$: the odd primes~$q_i$, at most
  $\mathrm{jacobiPrimes} = 16$, each with its reciprocal and, for odd~$e_i$, the table of the
  non-squares modulo~$q_i$, at most $\mathrm{jacobiTable} = 8192$ entries in all; and the sign for
  each $b^* \bmod 8$, the product of $J(-1 \mid b^*)$ for negative~$\ell$, $J(2 \mid b^*)^v$ and the
  signs of the reciprocity of the $q_i \equiv 3 \bmod 4$ with odd~$e_i$; none when a prime of~$\ell$
  lies beyond the table or a bound is passed. The test of~$b$ (\csource{find_points.c}{734}{759})
  divides~$2$ and the~$q_i$ out of~$b$ with the reciprocals and passes when the sign of the result
  times its entries modulo the~$q_i$ of odd exponent is~$1$. The contracts: tables exist only for
  $b_{\mathrm{high}} \le \mathrm{mulModLimit}$ and $\ell \ne 0$, and then describe~$\ell$ as above,
  with primes of the table and their reciprocals; for such tables and $1 \le b < 2^{64}$ with
  $b q < 2^{64}$ for every prime~$q$ of the tables, the test is the negated exclusive or of the sign
  and the entries at the part of~$b$ free of~$2$ and of these primes.
\end{definition}

The contract of the binary algorithm rests on invariants of its loops; $b^*$ is
\lname{Nat.coprimePart}, and the symbol is Mathlib's \verb|jacobiSym|.

\begin{lemma}[The Jacobi tests keep every point]\label{lem:jacobiTest}
  \uses{def:jacobi, def:binaryForm}
  \lean{Ratpoints.Sieve.run_jacobiTest, Ratpoints.Sieve.jacobiSym_eq_ite_foldl, Ratpoints.Sieve.jacobiSym_coprimePart_eq_one_of_sq_eq, Ratpoints.Sieve.run_jacobi_eq_one_of_sq_eq, Ratpoints.Sieve.run_jacobiTest_of_sq_eq}
  \leanok
  When the tables of~$\ell$ for~$b_{\mathrm{high}}$ exist, the test by tables passes
  $1 \le b \le b_{\mathrm{high}}$ exactly when $J(\ell \mid b^*) = 1$. Let~$d$ be even, $a$
  and~$y$ integers and~$b$ a natural number with $\gcd(a, b) = 1$ and $y^2 = F(a, b)$. Then
  $J(c_d \mid b^*) = 1$ for the largest divisor~$b^*$ of~$b$ prime to $2 c_d$; so for $b \ge 1$ the
  binary algorithm gives~$1$ on~$b$ and~$c_d$, and the test by tables of~$c_d$ for
  $b_{\mathrm{high}} \ge b$, when they exist, passes~$b$.
\end{lemma}

\begin{proof}
  \uses{def:jacobi, def:reciprocal, def:tableOfPrime, lem:exitByPrime}
  \leanok
  The denominators are below~$2^{32}$ and the primes of the table below~$2^{10}$, so the quotients
  by the reciprocals are exact (\lname{Ratpoints.Sieve.mulDiv\_magic}) and leave~$b^*$; by
  reciprocity each $J(q_i \mid b^*)$ of odd exponent is the Legendre symbol of $b^* \bmod q_i$,
  which the table gives, times a sign by $b^* \bmod 4$. For a triple of even degree, $c_d$ is a
  square modulo every prime of~$b$ prime to it
  (\lname{Ratpoints.Sieve.isSquare\_getD\_degree\_of\_sq\_eq\_binaryForm}).
\end{proof}

\section{The shapes of the loop}\label{sec:shapes}

\begin{definition}[The shapes of the loop over the denominators]\label{def:Shape}
  \uses{def:setupUs1, def:valuation, def:jacobi, def:Double, def:packing, def:bitArrays, def:exactBounds}
  \lean{Ratpoints.Sieve.Shape, Ratpoints.Sieve.runSetup, Ratpoints.Sieve.fst_run_runSetup, Ratpoints.Sieve.denInfoTest, Ratpoints.Sieve.wordsTest, Ratpoints.Sieve.Shape.dens, Ratpoints.Sieve.Shape.keep, Ratpoints.Sieve.wordsDens, Ratpoints.searchDens}
  \leanok
  The four shapes with their data (\csource{find_points.c}{3725}{3728},
  \csource{find_points.c}{3736}{3749}), which \verb|run_setup| chooses from the flags and returns
  with the bounds of the search intervals (\csource{find_points.c}{3778}{3806}), and their
  denominators from~$b_{\mathrm{low}}$ to~$b_{\mathrm{high}}$ in the order of the loop, each with
  its test of a single denominator:
  \begin{enumerate}
    \item under the flag of the squares, the squares~$k^2$ for
      $\lceil \sqrt{b_{\mathrm{low}}}\, \rceil \le k
      \le \lfloor \sqrt{b_{\mathrm{high}}}\, \rfloor$, with the class test;
    \item otherwise, when the set-up has returned the records and the divisors of
      Definition~\ref{def:setupUs1}, the divisors times squares: for each divisor~$s$ in turn,
      the~$s k^2$ of the range in increasing order of~$k$, with the class test and the valuation
      test, which~$b$ fails when the record of a prime excludes it
      (\csource{find_points.c}{3944}{3953});
    \item otherwise, under the flag of the checked loop, the words, with the forbidden divisors, the
      switch of the Jacobi test, its tables when the test is on and they exist, and~$c_d$: word by
      word, the $64 w + j$ of the range for the set bits~$j$, lowest first, of~$\mathrm{den\_bits}$
      and-ed with word $w \bmod p$ of each forbidden divisor~$p$ with words, tested by the tables,
      then by the forbidden divisors with masks, and, with the Jacobi test on and no tables, by the
      binary algorithm (\csource{find_points.c}{4013}{4041});
    \item otherwise every denominator, $b_{\mathrm{low}}, \dots, b_{\mathrm{high}}$, with the class
      test.
  \end{enumerate}
  The class test is that the pattern of the class of $b \bmod 64$ is not~$0$. The lists of the
  denominators of the shapes (\lname{Ratpoints.Sieve.Shape.dens}) and the choice of the test of a
  single denominator (\lname{Ratpoints.Sieve.Shape.keep}) are definitions of the proofs, the
  reference of the loops over the blocks, which test and sift each denominator as they reach it, as
  \verb|run_block| does; the tests themselves model the C code. The denominators of the search
  (\lname{Ratpoints.searchDens}) are those of the divisors times squares when~$d$ is odd,
  $|c_d| \ne 1$ and the set-up for~$b_{\mathrm{high}}$ returns divisors, and
  $b_{\mathrm{low}}, \dots, b_{\mathrm{high}}$ otherwise.
\end{definition}

The loops over the blocks (\lname{Ratpoints.Sieve.Shape.blockLoop}, Chapter~\ref{chap:moduli})
visit the denominators of the list of the shape in its order
(\lname{Ratpoints.Sieve.Shape.flatMap\_blockDens}). In the shape of the words they reduce the number
of the first word of a block modulo each forbidden divisor, with its reciprocal while that number is
at most~$\mathrm{mulModLimit}$ and by division otherwise, and step on from there
(\csource{find_points.c}{3981}{4006}).

\begin{lemma}[The shapes keep every point]\label{lem:shapes}
  \uses{def:Shape, def:twoAdic, def:setupUs1, def:forbidden, def:jacobi, def:binaryForm, def:packing}
  \lean{Ratpoints.Sieve.Shape.le_and_le_of_mem_dens, Ratpoints.Sieve.mem_squaresDens_iff, Ratpoints.Sieve.mem_squares1Dens_iff, Ratpoints.Sieve.nodup_squares1Dens, Ratpoints.Sieve.toList_run_wordsDens, Ratpoints.Sieve.run_denInfoTest_of_sq_eq, Ratpoints.Sieve.getLsbD_denBits_and_all_of_sq_eq, Ratpoints.Sieve.run_wordsTest_of_sq_eq, Ratpoints.Sieve.Shape.bits_ne_zero_of_mem_dens}
  \leanok
  Let $1 \le b_{\mathrm{low}}$. The denominators of every shape lie in the range: the squares are
  its squares in increasing order; for positive divisors, the divisors times squares are its~$s k^2$
  with~$s$ a divisor, without duplicates for distinct squarefree divisors; for forbidden divisors
  with positive primes and the words of \lname{Ratpoints.Sieve.forbiddenRow}, the words give, in
  increasing order, the denominators of the range with their bit in~$\mathrm{den\_bits}$ set that
  none of these primes divides. Let~$a$ and~$y$ be integers and $1 \le b \le H$ with
  $\gcd(a, b) = 1$ and $y^2 = F(a, b)$. For odd~$d$, $b$ passes the valuation test with the records
  of Definition~\ref{def:setupUs1}; for even~$d$ and admissible forbidden divisors for
  $b_{\mathrm{high}} = H$, the bit of $b \bmod 64$ in~$\mathrm{den\_bits}$ is set, no forbidden
  divisor with words divides~$b$, and~$b$ passes the tests of the words, with the tables for~$H$
  when the Jacobi test is on. When bit~$j$ of~$\mathrm{den\_bits}$ says whether the pattern of
  class~$j$ is not~$0$, a denominator of a shape that passes its test has a nonzero pattern.
\end{lemma}

\begin{proof}
  \uses{lem:twoAdic_point, lem:odd_denominators, lem:forbidden, lem:jacobiTest, def:bitArrays}
  \leanok
  The squares come from the integer square roots, $\lceil \sqrt{n}\, \rceil \le k$ exactly when
  $n \le k^2$; $s k^2$ determines the squarefree~$s$ and~$k$
  (\lname{Nat.eq\_of\_mul\_sq\_eq\_mul\_sq}); and bit~$t$ of word $w \bmod p$ of a forbidden divisor
  is clear exactly when $p \mid 64 w + t$
  (\lname{Ratpoints.Sieve.getLsbD\_getElem!\_forbiddenRow}). The tests pass~$b$ by
  Lemmas~\ref{lem:twoAdic_point}, \ref{lem:odd_denominators}, \ref{lem:forbidden}
  and~\ref{lem:jacobiTest}.
\end{proof}

\section{The order of the denominators}\label{sec:searchDens}

The shapes visit their denominators in increasing order, except the divisors times squares, which
go divisor by divisor. The denominators that the other shapes and the tests leave out carry no point
(Lemma~\ref{lem:shapes}), so among those that carry points only the divisor shape changes the
order. For $y^2 = 1 + 2x$ and $H = 20$, with the divisors~$1$ and~$2$, the loop visits
$1, 4, 9, 16$ and then $2, 8, 18$ ($4$ and~$16$ fail the valuation test of the prime~$2$), so that
\ratpoints{} prints $(-4 : 3 : 9)$ before $(-1 : 0 : 2)$.

\begin{lemma}[The denominators of the search]\label{lem:searchDens}
  \uses{def:setupUs1, def:Shape, def:orderedPoints, def:bruteForce, def:maxDegree}
  \lean{Ratpoints.searchDens_one_eq_searchDenominators, Ratpoints.searchDens_eq_filter, Ratpoints.perm_filter_bruteForce_append_flatMap_searchDens}
  \leanok
  For $b_{\mathrm{high}} \ge 1$ the denominators of the search of~$c$ of its degree from~$1$
  to~$b_{\mathrm{high}}$ are
  $\mathrm{searchDenominators}(\mathrm{maxDivisors}, c, b_{\mathrm{high}})$
  (Definition~\ref{def:orderedPoints}), and those from $b_{\mathrm{low}} \ge 1$ are those from~$1$
  that are at least~$b_{\mathrm{low}}$. For $1 \le b_{\mathrm{low}}$ and
  $1 \le b_{\mathrm{high}} \le H$, the points that the brute-force reference finds over $(1, 0)$,
  the points at infinity, followed by those over the denominators of the search with the numerators
  $-H, \dots, H$, are a permutation of the points of $\mathrm{bruteForce}(c, H)$ with $z = 0$ or
  $b_{\mathrm{low}} \le z \le b_{\mathrm{high}}$.
\end{lemma}

\begin{proof}
  \uses{def:setupUs1, lem:odd_denominators, lem:shapes}
  \leanok
  The specification takes the divisors times squares under the same condition, with the divisors
  $\mathrm{binaryDivisors}$ (Definition~\ref{def:setupUs1}); from~$1$ the divisor shape lists the
  $s k^2 \le b_{\mathrm{high}}$ divisor by divisor
  (\lname{Ratpoints.Sieve.squares1Dens\_one\_eq\_filter}). The divisor shape holds every
  denominator of a point in the range once (Lemmas~\ref{lem:odd_denominators}
  and~\ref{lem:shapes}), so its points are a permutation of those of the whole range
  (\lname{Ratpoints.perm\_flatMap\_squares1Dens}).
\end{proof}

With the block lemma (Theorem~\ref{thm:blocks}), these lemmas give the list that a search delivers
(Chapter~\ref{chap:theorems}), on which the order theorem (Theorem~\ref{thm:order}) and the
membership theorem (Theorem~\ref{thm:membership}) rest.

\section{Where the model differs from the C code}\label{sec:denominatorsDeviations}

One function serves for \verb|valuation| and \verb|valuation1|, one for \verb|jacobi| and
\verb|jacobi1|; the square shapes take exact integer square roots, which agree with the C code's
corrected square roots in \verb|double| on the arguments it passes; and the symbol
$J(c_d \mid p) = -1$ of the further primes is read off the table of the squares modulo~$p$. The
masks of the forbidden valuations have no bound $m < 63$, which never binds for a \verb|long|
$b_{\mathrm{high}}$ (\lname{Ratpoints.Sieve.lt\_forty\_of\_pow\_le}), so that the model is sound for
every height. The other deviations, natural numbers for the C integer types and no failure for
lack of memory, are those of the whole model (Chapter~\ref{chap:ccode}).
